\section{Ex ante or ex post? The adoption absorber}
\label{sec:exante}

The instruments in Section~\ref{sec:shield} are all \emph{ex post}: deployed once
the shock has landed. Because the adoption margin responds to the same relative
price the crisis moves, a natural alternative is to act \emph{ex ante}, to lean
on the margin in the steady state with a permanent carbon tax $\tau_b$ on brown
energy, leaving the economy greener and, one might hope, less exposed. This
section prices that hope. The answer is that ex-ante greening is a poor substitute
for ex-post insurance, and the reason is intrinsic to the channel this paper adds:
the adoption margin is an endogenous shock absorber, and pre-greening spends it.

\subsection{A welfare-optimal permanent carbon tax}
\label{sec:carbontax}

Holding $\psi_g$ fixed at its no-tax value and letting $D^{G}_{ss}$ float, a
carbon tax raises the effective brown/green price ratio and lifts steady-state
adoption: $D^{G}_{ss}$ moves from $5\%$ at $\tau_b=0$ to $8.9\%$ at $\tau_b=0.10$
and $20.1\%$ at $\tau_b=0.30$ (Figure~\ref{fig:carbon_optimal},
right). Steady-state welfare is hump-shaped in $\tau_b$
(Figure~\ref{fig:carbon_optimal}, left), rising to a peak of $+0.043\%$
at $\tau_b^{*}\approx 10\%$ and turning negative beyond
$\tau_b\approx 25\%$. The model contains \emph{no} climate
externality, so this interior optimum is not Pigouvian. It is a second-best
correction of the adoption friction: green energy is structurally cheaper
($\bar P^{E}_{G}=0.8<\bar P^{E}_{B}=1$) but the logit taste scale
$\sigma_\varepsilon$ and the switching cost $\psi_g$ leave green under-adopted, so
a modest tax raises welfare until the tax distortion dominates. We treat this as a
diagnostic of the channel, not a policy recommendation: it would be swamped by
any explicitly modelled externality or abatement cost. The carbon revenue is
rebated lump-sum here; Appendix~\ref{sec:recycle} studies the budget-neutral
alternative that earmarks it to a green switching subsidy, which greens the steady
state \emph{more} yet leaves a residual brown pool \emph{less} responsive to the
crisis: the same absorber logic as Section~\ref{sec:routeA}, across instruments.

\begin{figure}[H]\centering
\includegraphics[width=\textwidth]{output/fig_carbon_optimal_import.pdf}
\caption{The welfare-optimal permanent carbon tax. Left: standing CEV of the
$\tau_b$ steady state relative to the no-tax baseline (no crisis), hump-shaped
with an interior optimum near $\tau_b^{*}\approx10\%$. Right: the induced
steady-state green share $D^{G}_{ss}$. $\psi_g$ is held fixed; the carbon revenue
is rebated lump-sum.}
\label{fig:carbon_optimal}
\end{figure}

\subsection{Ex-ante greening barely insures}
\label{sec:versionA}

We now hit two economies with the \emph{same} price shock: the baseline
($D^{G}_{ss}=5\%$) and a ``prepared'' economy carrying the welfare-optimal
$\tau_b=0.10$ ($D^{G}_{ss}=8.9\%$). The prepared economy is essentially no better
insured: its crisis CEV is $-1.85\%$ against $-1.87\%$ for laissez-faire, a
difference of two basis points, and both are far worse than either ex-post
instrument (cap $-1.15\%$, transfer $-0.82\%$; Table~\ref{tab:exante_expost_import}).
The greening the tax buys is real, $D^{G}_{ss}$ is $78\%$ higher, but too small a
share of the economy to move aggregate exposure, and the tax's crisis-time value
is a rounding error next to a transfer that puts cash in high-MPC hands exactly
when it binds. Consistent with Section~\ref{sec:cap}, the combination ``ETS $+$
cap'' tracks the cap alone: once the price is frozen the adoption margin is shut
and the greener starting point is irrelevant.

The two economies are worth seeing side by side. Figure~\ref{fig:ets_irf}
overlays their responses to the same brown-price shock. The aggregates, output,
consumption, inflation, net exports, are nearly indistinguishable; the two differ
visibly only in brown energy use, the green share $\DG$, the switching flow, and
the brown price. The greener starting point changes the \emph{composition} of the
response, not its aggregate size.

\begin{figure}[H]\centering
\IfFileExists{output/fig_ets_irf_price_import.pdf}{\includegraphics[width=\textwidth]{output/fig_ets_irf_price_import.pdf}}{\textit{[figure: run run\_ets\_cross\_section.py]}}
\caption{The same brown-price shock hitting the baseline ($\DG_{ss}=5\%$) and the
prepared ETS ($\tau_b=0.10$, $\DG_{ss}=8.9\%$) steady states, in level deviations
from each economy's own steady state. Aggregates coincide; brown energy, the
green share, switching and the brown price differ.}
\label{fig:ets_irf}
\end{figure}

Figure~\ref{fig:cross_section} and Table~\ref{tab:cross_section_import} decompose
the consumption response by the household's current technology. Switchers consume
roughly twice the average brown household in the steady state, so a naive
current-group split is contaminated by selection; we use the common-steady-state
counterfactual to hold group populations fixed, read the clean per-capita response
of each group off the frozen path, and report the adoption/migration term (full
minus frozen) separately. Same aggregate loss, different cross-section: under the
ETS both brown and green incumbents lose \emph{less} (per-capita $-92\%$ against
$-97\%$, and $-28\%$ against $-37\%$ of their own steady state over $H=24$), and
that relief is exactly offset by a heavier adoption term ($-22.8$ against $-13.2$,
cumulative $\times 100$), more households paying $\psi_g$ to switch during the
crisis. Green incumbents, fully insulated from the brown price, still
lose about $1.5\%$ per quarter, through the general-equilibrium income channel
rather than energy.

\begin{figure}[H]\centering
\IfFileExists{output/fig_cross_section_price_import.pdf}{\includegraphics[width=\textwidth]{output/fig_cross_section_price_import.pdf}}{\textit{[figure: run run\_ets\_cross\_section.py]}}
\caption{Consumption by technology group under the brown-price shock, baseline
and ETS. Left: per-capita consumption of brown and green incumbents at fixed
populations (\% of own steady state). Centre and right: contributions to $dC$
(fixed-population brown, fixed-population green, and the adoption/migration term).}
\label{fig:cross_section}
\end{figure}

\input{output/tab_cross_section_import.tex}

\input{output/tab_exante_expost_import.tex}

\subsection{The green subsidy at the trough}
\label{sec:greensub}

The green/adoption subsidy is worth studying as a crisis instrument in its own
right (Figure~\ref{fig:green_subsidy}): the government pays a fraction of the
switching cost during the acute crisis. It roughly doubles the adoption response
(peak $D^{G}$ of $18.0$ vs.\ $8.9$ pp; switchers $2.3$ vs.\ $1.0$ pp) at trivial
fiscal cost ($2.0$ against the cap's $53.9$), yet its CEV is the worst in
Table~\ref{tab:summary_price_import}, $-2.03\%$, below laissez-faire, and it
deepens the output loss ($\sum y=-70.1$). Paying households to incur the real
switching cost $\psi_g D^{sw}$, booked as an import (Section~\ref{sec:bop}),in
the middle of a consumption crisis destroys welfare even as it greens the economy.
The adoption margin is valuable, but not as something to force at the trough.

\begin{figure}[H]\centering
\IfFileExists{output/fig_green_subsidy_import.pdf}{\includegraphics[width=\textwidth]{output/fig_green_subsidy_import.pdf}}{\textit{[figure: run run\_green\_subsidy.py]}}
\caption{The green/adoption subsidy at the trough, against laissez-faire, the
price cap and the Slutsky transfer (brown-price shock). It doubles the adoption
response ($D^{G}$, $D^{\mathrm{sw}}$) at trivial fiscal cost, yet forcing the real
switching cost in the middle of the consumption crisis leaves output and private
welfare below laissez-faire.}
\label{fig:green_subsidy}
\end{figure}

\subsection{Persistence exhausts the absorber}
\label{sec:routeA}

The obvious defence of the ex-ante case is persistence: a longer-lived shock
should let the greener economy's lower exposure compound over more periods.
Figure~\ref{fig:persistence} and Table~\ref{tab:persistence_import} test it by
sweeping the shock half-life and re-solving only the impulse response (the steady
state does not move with persistence), decomposing the ex-ante number into its
constant standing part and a \emph{crisis-only greening dividend}, defined as the
prepared economy's crisis CEV minus the standing benefit minus laissez-faire.

The defence fails, and instructively, though the failure is one of welfare, not
of switching volume. The crisis-only greening dividend does not grow with
persistence; if anything it falls with the half-life of the shock. Its level is at
most a few basis points at every horizon, inside the first-order approximation
error of the CEV (Appendix~\ref{sec:nlcev}), so we do not sign it: it is zero to
first order. The ex-post cap is different in kind, its crisis protection rises
monotonically, from $+0.24\%$ to $+1.36\%$ as a longer freeze shields more of the
transition, an order of magnitude larger than the approximation error and robust
nonlinearly.

The reason is the adoption margin itself. A persistent brown-price shock makes
switching attractive for many quarters, so the \emph{browner} laissez-faire
economy adopts heavily \emph{during} the crisis and reaps that windfall over the
whole episode, the margin is an endogenous shock absorber. The prepared economy
has already spent much of that absorber in the steady state. It is not that it
adopts less in the crisis, with a standing tax it in fact adopts \emph{more}
(peak $\DG$ of $15.0$ against $8.9$~pp; Section~\ref{sec:versionA}), but that the
marginal crisis switch is worth less to it: the households the tax has already
moved to green are inframarginal to the shock, and the extra crisis adoption pays
the real cost $\psi_g$ at the trough (Section~\ref{sec:greensub}). The
welfare-relevant absorber is the \emph{option} to green cheaply exactly when the
crisis makes it optimal, and the browner economy enters with that option more
fully intact. Pre-greening raises the steady-state green share and even the crisis
switching volume; what it does not raise is the crisis \emph{welfare} value of the
margin. This is the sharpest form of the paper's
thesis. The instruments that freeze the adoption technology, the cap here, and by
construction the frozen-technology models of \citet{Bayer2026},
\citet{Langot2026} and \citet{ARS2024}, miss an absorber that works best when
left loaded, i.e.\ when the economy enters the crisis with room to green rather
than already green.

\begin{figure}[H]\centering
\IfFileExists{output/fig_persistence_import.pdf}{\includegraphics[width=\textwidth]{output/fig_persistence_import.pdf}}{\textit{[figure output/fig\_persistence\_import.pdf: run run\_persistence.py]}}
\caption{Persistence and the two policies. Left: total CEV of laissez-faire, the
ex-post cap and the ex-ante ETS as the shock half-life rises. Right: the
crisis-only gain over laissez-faire, the cap's protection grows with persistence,
while the ex-ante \emph{greening dividend} does not grow, it is zero to first
order, because a more persistent shock lets the browner economy exploit the
adoption absorber the prepared economy has already spent.}
\label{fig:persistence}
\end{figure}

\IfFileExists{output/tab_persistence_import.tex}{\input{output/tab_persistence_import.tex}}{\textit{[table \texttt{tab\_persistence\_import}: run \texttt{run\_persistence.py}]}}

\paragraph{Takeaway.} Three results point the same way. Energy-indexed targeting
is dominated by a flat transfer of equal budget (Section~\ref{sec:flat}); a
permanent carbon tax is mildly welfare-improving in the steady state through a
non-Pigouvian adoption correction, but essentially worthless as crisis insurance
(Sections~\ref{sec:carbontax}--\ref{sec:versionA}); and its insurance value does not
grow, it is zero to first order, as shocks become persistent, because greening
ex ante exhausts the adoption absorber (Section~\ref{sec:routeA}). The margin this paper adds is not a
lever a planner should pre-pull; it is a stabiliser that pays off precisely when
the economy is left free to adopt in the crisis.
